\chapter{Sprint 4 : Nutrition, Communication et Assistance}
\minitoc
\newpage

\section{Sprint IV : Nutrition, Communication et Assistance}

\subsection{Objectif du Sprint}
La performance sportive ne peut être dissociée d'une nutrition adaptée ni d'une communication régulière entre l'athlète et le professionnel. Ce quatrième sprint regroupe les fonctionnalités de finalisation métier : module nutritionnel, suivi des objectifs, messagerie instantanée, sessions de visioconférence et assistance IA.

L'approche adoptée est centrée sur les données et l'accompagnement : chaque repas consommé est enregistré avec ses macronutriments, tandis que les échanges, notifications et sessions permettent de maintenir un suivi continu.

\subsection{Backlog du Sprint IV}

\begin{table}[H]
\centering
\renewcommand{\arraystretch}{1.4}
\resizebox{\textwidth}{!}{%
\begin{tabular}{|c|c|p{5cm}|p{6.5cm}|c|}
\hline
\rowcolor[HTML]{DAE8FC}
\textbf{ID} & \textbf{Acteur} & \textbf{User Story} & \textbf{Tâche Technique} & \textbf{Priorité} \\ \hline
4.1 & Coach/Nutritionniste & Créer un plan alimentaire (Diet Builder) & Entités DietPlan/DietDay/Meal & Must \\ \hline
4.2 & Coach/Nutritionniste & Assigner un plan à un athlète & Assignation contrôlée du plan alimentaire & Must \\ \hline
4.3 & Athlète & Logger un repas consommé & Entité MealLog et journal alimentaire & Must \\ \hline
4.4 & Athlète & Consulter le bilan macro du jour & Calcul du bilan macro journalier & Must \\ \hline
4.5 & Athlète & Définir ses objectifs caloriques personnels & Entité Goal et validation des objectifs & Should \\ \hline
4.6 & Coach & Consulter la conformité nutrition d'un athlète & Indicateurs de conformité nutritionnelle & Should \\ \hline
4.7 & Athlète & Voir l'historique des journées passées & Historique nutritionnel & Could \\ \hline
4.8 & Tous & Envoyer et recevoir des messages en temps réel & Socket.io, événements send/receive, entité Message & Should \\ \hline
4.9 & Tous & Réserver une session de visioconférence & Génération de lien Jitsi et calendrier & Could \\ \hline
4.10 & Coach/Athlète & Consulter un copilote IA pour des conseils & Intégration Gemini API & Could \\ \hline
\end{tabular}%
}
\caption{Backlog du Sprint 4 -- Nutrition, Communication et Assistance}
\end{table}

\section{Analyse}

\subsection{Diagramme de cas d'utilisation du quatrième sprint}
\begin{figure}[H]
\centering
\includegraphics[width=0.88\textwidth]{Sprint4_Cas_Utilisation.png}
\caption{Diagramme de cas d'utilisation -- Sprint 4}
\end{figure}

\subsection{Tableau descriptif du cas d'utilisation « ... »}

\subsection{Tableau descriptif du cas d'utilisation « ... »}

\subsection{Tableau descriptif du cas d'utilisation « ... »}

\subsection{Tableau descriptif du cas d'utilisation « ... »}

\subsection{Tableau descriptif du cas d'utilisation « ... »}

\subsection{Tableau descriptif du cas d'utilisation « ... »}

\subsection{Tableau descriptif du cas d'utilisation « ... »}

\subsection{Tableau descriptif du cas d'utilisation « ... »}

\section{Conception}

\subsection{Diagrammes de séquences}

\subsection{A) Diagramme de séquence « ... »}
\begin{figure}[H]
\centering
\includegraphics[width=0.85\textwidth]{Sprint4_Sequence_LogMeal.png}
\caption{Diagramme de séquence -- Enregistrement d'un repas consommé}
\end{figure}

\subsection{B) Diagramme de séquence « ... »}

\subsection{C) Diagramme de séquence « ... »}

\subsection{D) Diagramme de séquence « ... »}

\subsection{Diagramme de classes du quatrième sprint}
\begin{figure}[H]
\centering
\includegraphics[width=0.9\textwidth]{Sprint4_Classes.png}
\caption{Diagramme de classes -- Plans nutritionnels et logs}
\end{figure}

La structure de données nutritionnelle est symétrique à celle des programmes sportifs :
\begin{itemize}
    \item \textbf{DietPlan :} Conteneur principal (nom, objectif calorique global).
    \item \textbf{DietDay :} Un jour de la semaine du plan avec ses macros cibles (protéines, glucides, lipides).
    \item \textbf{Meal :} Un repas précis du jour (petit-déjeuner, déjeuner, collation, dîner) avec le détail des apports nutritionnels.
    \item \textbf{MealLog :} Enregistrement réel de ce que l'athlète a consommé ce jour-là.
\end{itemize}

\section{Réalisation}

\subsection{Description des interfaces}

\section*{Conclusion du Sprint IV}
\addcontentsline{toc}{section}{Conclusion du Sprint IV}
\mtcaddsection
\markboth{Conclusion du Sprint IV}{}
Le sprint finalise les fonctionnalités métier principales de GOSPORT. Les nutritionnistes et coachs peuvent concevoir des plans diététiques détaillés, les athlètes peuvent suivre leurs repas et objectifs, et les échanges sont renforcés par la messagerie, les sessions de visioconférence et l'assistance IA. Le projet dispose ainsi d'un parcours complet : inscription, mise en relation, entraînement, nutrition et accompagnement continu.
